\Askhsh{31680_SOLUTION}{1}
\begin{ylh}{1}
    \item [Διαφορικός Λογισμός]
\end{ylh}
Α) Είναι $L = AB + B\Gamma$. Από την εφαρμογή του πυθαγορείου θεωρήματος έχουμε ότι:
\[
    AB^2 = \Delta B^2 + \Delta A^2 = x^2 + (1{,}5)^2 = x^2+2{,}25
    \quad \text{και} \quad
    B\Gamma^2 = BE^2 + E\Gamma^2 = (1{,}7-x)^2 + \alpha^2.
\]
Άρα είναι
\[
    AB = \sqrt{x^2+2{,}25} \quad \text{και} \quad B\Gamma = \sqrt{(1{,}7-x)^2+\alpha^2}
    \quad \text{και} \quad L(x)=\sqrt{x^2+2{,}25}+\sqrt{(1{,}7-x)^2+\alpha^2}.
\]
\begin{alist}
\item [β)] Είναι
\[
    L'(x)=\frac{x}{\sqrt{x^2+2{,}25}} - \frac{1{,}7-x}{\sqrt{(1{,}7-x)^2+\alpha^2}}, \quad x \in \left(0, \frac{17}{10}\right).
\]
\begin{rlist}
\item Γνωρίζουμε ότι το $L$ ελάχιστο όταν το $B$ απέχει $1{,}02$ μέτρα από το $\Delta$. Επειδή $1{,}02 \in \left(0, \frac{17}{10}\right)$, η συνάρτηση $L$ είναι παραγωγίσιμη στο $1{,}02$ και το $L(1{,}02)$ είναι τοπικό ακρότατο (ως ελάχιστη τιμή) από το θεώρημα Fermat έπεται ότι $L'(1{,}02)=0$.
\item Επειδή $L''(x)>0$ για κάθε $x \in \left(0, \frac{17}{10}\right)$, η συνάρτηση $L'$ είναι γνησίως αύξουσα.
Ακόμη η $L'$ είναι συνεχής ως παραγωγίσιμη και από το θεώρημα Fermat είναι $L'(1{,}02)=0$, οπότε:
\end{rlist}
\end{alist}